\documentclass[12pt,a4paper]{article}
\usepackage[russian]{babel}
\usepackage{fontspec}
\usepackage{geometry}
\usepackage{booktabs}
\usepackage{longtable}
\usepackage{array}
\usepackage{multirow}
\usepackage{enumitem}
\usepackage{titlesec}
\usepackage{xcolor}
\usepackage{siunitx}

\usepackage{pdfpages}

\usepackage{tablefootnote}

% Дополнительные пакеты для данного документа
\usepackage{caption}
\usepackage{subcaption}
\usepackage{listings}
% \usepackage{hyperref}

\geometry{a4paper, margin=25mm}
\setmainfont{Liberation Serif}
\setsansfont{Liberation Sans}

\definecolor{adblue}{RGB}{0,84,166}

\titleformat{\section}{\large\bfseries\color{adblue}}{}{0em}{}[\titlerule]
\titleformat{\subsection}{\normalsize\bfseries}{}{0em}{}
\titleformat{\subsubsection}{\normalsize\itshape}{}{0em}{}

\sisetup{output-decimal-marker={,}}

\usepackage{tikz}
\usetikzlibrary{positioning, calc, backgrounds, math, 3d, perspective, arrows.meta, graphs, graphdrawing, fit, shapes.geometric, trees}
    

\usepackage{pgfplots} % для построения графиков
\pgfplotsset{compat=1.18} % версия для совместимости
\usepackage{tkz-euclide}

% --- Подключение пакета алгоритмов (без флага algo2e, чтобы вернуть \begin{algorithm}) ---
\usepackage[ruled,vlined]{algorithm2e}
%\usepackage{caption}
\newcommand{\Gray}[1]{\textcolor{gray}{#1}} 
% Настройка подписей
\DeclareCaptionLabelFormat{gost}{#1 #2}
\captionsetup[table]{labelformat=gost,labelsep=endash,justification=justified,singlelinecheck=false,font=normal}
\captionsetup[figure]{labelformat=gost,labelsep=endash,justification=centering,font=normal}

% Локализация algorithm2e на русский (ОБЯЗАТЕЛЬНО после загрузки пакета)
\SetAlgorithmName{Алгоритм}{}{}
\SetAlgoCaptionSeparator{---}
\SetKwInput{Input}{Вход}
\SetKwInput{Output}{Выход}

\usepackage{url}
% Говорим пакету использовать моноширинный шрифт (как у texttt)
\Urlmuskip=0mu plus 1mu % гибкие пробелы для формул, если нужно
\DeclareUrlCommand\code{\urlstyle{tt}} 

% Библиотека схемотехники
\usepackage[siunitx, RPvoltages, american]{circuitikzgit}
\ctikzset{
    amplifiers/fill=cyan,
    sources/fill=green!20,
    diodes/fill=white,
    resistors/fill=violet,
    grounds/scale=1.2,
}

\usepackage{iftex}

\ifXeTeX % Если используется TeTeX или TeLaTeX
  \typeout{Режим XeTeX}
  \usepackage{xltxtra}
  \usepackage{fontspec}
  \usepackage{polyglossia}
  \setmainlanguage{russian}
  \setotherlanguage{english}
  \setkeys{russian}{babelshorthands=true} % По идее включает <<, >>, -- ---

  \setmainfont{Times New Roman}[Ligatures=TeX]
  \setromanfont{Times New Roman}[Ligatures=TeX]
  \setsansfont{Arial}[Ligatures=TeX]
  \setmonofont{Courier New}[Ligatures=TeX] 

  \newfontfamily{\cyrillicfont}{Times New Roman}
  \newfontfamily{\cyrillicfontrm}{Times New Roman}
  \newfontfamily{\cyrillicfonttt}{Courier New}
  \newfontfamily{\cyrillicfontsf}{Arial}
\fi

\ifLuaTeX % Если используется LuaLatex
  \typeout{Режим LuaTeX}
%  \usepackage{luatextra}
  \usepackage{fontspec}
  \defaultfontfeatures{Ligatures={TeX}}
  \usepackage{polyglossia}
  \setmainlanguage{russian}
  \setotherlanguage{english}
  \setmainfont{Times New Roman}[Script=Cyrillic, Ligatures=TeX]
  \setromanfont{Times New Roman}[Script=Cyrillic, Ligatures=TeX]
  \setsansfont{Arial}[Script=Cyrillic, Ligatures=TeX]
  \setmonofont{Courier New}[Ligatures=TeX]
  \newfontfamily{\cyrillicfont}{Times New Roman}[Script=Cyrillic, Ligatures=TeX]
  \newfontfamily{\cyrillicfontrm}{Times New Roman}[Script=Cyrillic, Ligatures=TeX]
  \newfontfamily{\cyrillicfonttt}{Courier New}
  \newfontfamily{\cyrillicfontsf}{Arial}[Script=Cyrillic, Ligatures=TeX]

  \usegdlibrary{trees}
  \usegdlibrary{force}
\fi

\ifPDFTeX % Если используется 8битный PDFTex
  \typeout{8-битный режим pdfTex}
  % For pdfTeX we must use old
  % 8-bit font technologies
  \usepackage[T2A]{fontenc}
  \usepackage[english,russian]{babel}
  \usepackage[utf8]{inputenc}
  \usepackage[english, russian]{babel}

  %\usepackage{newtxmath}
  \usepackage{amsmath}
  \usepackage{amsthm}
  \usepackage{amssymb}
  \usepackage{amsfonts}
  \usepackage{mathtext}
\fi

\usepackage{amsthm}  % Возможности AMSMath
\usepackage{amsmath}
\usepackage{amssymb}

\usepackage{esvect} % Для красивых стрелок
\usepackage{bm} % Для жирных символов

%Hyphenation rules
\usepackage{hyphenat}
\hyphenation{ма-те-ма-ти-ка вос-ста-нав-ли-вать ото-бра-жа-ют-ся сге-не-ри-ро-ван-ном сим-во-лы}
% \usepackage[english, russian]{babel}

\usepackage{graphicx}
\graphicspath{ {./images/} }

\usepackage{empheq}

\usepackage{float}

% \usepackage{comment}

\usepackage[hidelinks,
    unicode=true,          % поддерживает Unicode в закладках
    psdextra=true,      % автоматически конвертирует простую математику в текст
    % focusdefault=false
]{hyperref}


% Окружение "Теорема"
\newtheorem{theorem}{Теорема}[section]
\newtheorem{corollary}{Corollary}[theorem]
\newtheorem{lemma}[theorem]{Lemma}

% Окружение "Определение"
\theoremstyle{definition} % Прямой шрифт, нумерация
\newtheorem{definition}{Определение}[section]
\newtheorem{example}{Пример}[section]

% Окружение "Замечание"
\theoremstyle{remark} % Прямой шрифт, без нумерации
% Окуржение "Решение"
\newtheorem*{solution}{Решение}
\newtheorem*{remark}{Замечание}

\usepackage{polynom}
\usepackage{tabularx}
\usepackage{cancel}

\addto\captionsrussian{%
  \renewcommand{\figurename}{Рис.}%
  \renewcommand{\tablename}{Табл.}%
}

\renewcommand{\arraystretch}{1.5} % Высота строки таблицы

\renewcommand{\labelitemi}{\(\circ\)} % Бюллет -- пустой кружок в itemize


\graphicspath{{figures/} {images/}}

\begin{document}


\subsection{Теоретическое обоснование спектрального анализа функции потерь в задачах демозаики}

Классические пространственные критерии оптимизации, такие как средний абсолютный шаг ошибки ($L_1$-loss) и среднеквадратичное отклонение ($L_2$-loss), минимизируют попиксельную разность амплитуд в пространственной области $\mathbb{R}^{2}$. Однако в задачах восстановления сырых RAW-данных (дебаеризации) указанные метрики демонстрируют принципиальную ограниченность, обусловленную эффектом спектрального наложения (aliasing) и высокочастотной пульсацией периодической структуры мозаики фильтров Байера (RGGB).

Применение одномерных или двумерных пространственных сверток приводит к сглаживанию (мылу) микроконтуров и потере хроматической дискретизации на экстремальных точках спектра (например, в пурпурных и фиолетовых диапазонах). Для устранения этого эффекта в рамках данного исследования математический аппарат минимизации ошибки переведён в частотную область посредством дискретного двумерного преобразования Фурье ($\text{2D-DFT}$).

Переход от пространственных координат к спектральным плотностям для предсказанного тензора $\hat{Y}$ и эталонного высококонтрастного изображения $Y$ осуществляется на основе классического преобразования:

\begin{equation}
\mathcal{F}(X)_{u,v} = \sum_{x=0}^{M-1} \sum_{y=0}^{N-1} X_{x,y} \cdot \exp\left( -j 2\pi \left( \frac{ux}{M} + \frac{vy}{N} \right) \right)
\end{equation}

где $M$ и $N$ — пространственные размеры обрабатываемого патча изображения. Для стабилизации фазовых сдвигов и централизации низкочастотных компонент (постоянной составляющей яркости) к центру спектра применяется оператор циклического сдвига знаков четвертей комплексной плоскости:

\begin{equation}
\mathcal{F}_{\text{shift}}(X) = \mathcal{Shift}\left( \mathcal{F}(X) \right)
\end{equation}

Амплитудный спектр частот, определяемый как модуль комплексного числа $| \mathcal{F}_{\text{shift}}(X) | = \sqrt{\text{Re}^2 + \text{Im}^2}$, характеризуется колоссальным динамическим диапазоном (энергия низких частот на порядки превосходит энергию высоких частот контуров). Прямая минимизация среднеквадратичного расстояния на таких массивах приводит к тому, что нейросеть оптимизирует только общую освещенность кадра, игнорируя мелкоструктурный байеровский муар.

Для компрессии динамического диапазона частот и выравнивания градиентных потоков весов автодифференцирования введена нелинейная логарифмическая шкала с масштабным коэффициентом сжатия $\gamma = 100.0$:

\begin{equation}
\mathcal{A}_{\log}(X) = \frac{\ln\left(1 + \gamma \cdot |\mathcal{F}_{\text{shift}}(X)|\right)}{\ln(1 + \gamma)}
\end{equation}

Знаменатель $\ln(1 + \gamma)$ является константным нормирующим делителем, переводящим амплитуды в шкалу $[0, 1]$. Дифференциал логарифмической функции потерь $L_{\text{FFL\_Log}}$ по предсказанным амплитудам $\mathcal{A}_{\log}(\hat{Y})$ формирует устойчивый спектральный фильтр, который наказывает нейросетевую модель за малейшие отклонения в периодических высоких частотах. Это вынуждает глубокие блоки архитектуры NAFNet выучивать честное аналитическое распределение цвета, полностью подавляя паразитную шахматную сетку деградации сенсора.

\section{Алгоритм}
\begin{algorithm}[H]
\DontPrintSemicolon
\SetKwInOut{KwIn}{Вход}
\SetKwInOut{KwOut}{Выход}

\TitleOfAlgo{Сквозной конвейер восстановления RAW-данных: Модуль оркестрации}

\KwIn{Каталог $D_{\text{source}}$, конфигурация \texttt{config.yml}, масштабы патчей $S_{\text{lq}}, S_{\text{gt}}$, масштаб $F_{\text{scale}}$.}
\KwOut{Сгенерированный фазово-выровненный датасет в $D_{\text{root}}$, диагностические превью $\text{Img}_{\text{lq}}, \text{Img}_{\text{gt}}$.}

\BlankLine
\tcc*[l]{Этап I: Инициализация и верификация подсистемы ввода-вывода}
Парсинг и валидация аргументов командной строки через модуль \texttt{config.py}\;
Инициализация двухпоточного логгера \texttt{setup\_logger()} в файл \texttt{pipeline.log}\;
\If{выставлен флаг \texttt{--clean-training}}{
    Выборочное удаление старых каталогов \texttt{checkpoints}, \texttt{tb\_logs}, \texttt{val\_predictions}\;
}

\BlankLine
\tcc*[l]{Этап II: Синхронная генерация синтетического датасета}
\If{каталог датасета пуст {\bf или} выставлен флаг \texttt{--clean-dataset}}{
    Очистка директорий \texttt{train} и \texttt{test} внутри корня $D_{\text{root}}$\;
    Разбиение файлов из $D_{\text{source}}$ на подмножества $\text{Files}_{\text{train}}$ и $\text{Files}_{\text{test}}$ в пропорции $\text{ratio}$\;
    \ForEach{изображение $I_{\text{rgb}}$ в $\text{Files}_{\text{train}}, \text{Files}_{\text{test}}$}{
        Чтение матрицы пикселей и приведение к 8-битному RGB-пространству\;
        Вырезание центрального патча $I_{\text{crop}}$ размера $S_{\text{gt}} \times S_{\text{gt}}$ с соблюдением чётности Bayer-фазы\;
        Применение PSF-размытия оптической системы: $I_{\text{blur}} = \mathcal{Filter2D}(I_{\text{crop}}, \text{Kernel}_{\text{PSF}})$\;
        Накладывание маски мозаики фильтров: $I_{\text{bayer}} = \mathcal{ApplyBayerMask}(I_{\text{blur}}, \text{'RGGB'})$\;
        Генерация аддитивно-мультипликативного шума на основе SNR в дБ\;
        Упаковка 2D-мозаики в 4-канальный массив субканалов: $I_{\text{packed}} = \mathcal{ExtractSubchannels}(I_{\text{bayer}})$\;
        Сохранение пар: $I_{\text{packed}} \rightarrow \text{uint16 tiff}$, $I_{\text{crop}} \rightarrow \text{uint8 png}$\;
    }
}

\BlankLine
\tcc*[l]{Этап III: Диагностический визуальный контроль данных}
Извлечение случайного верификационного сэмпла из класса \texttt{CustomNEFPairDataset}\;
Применение быстрой билинейной интерполяции к 4-канальному входу\;
Конвертация цветового пространства из RGB в BGR: $I_{\text{bgr}} = \mathcal{Cv2Color}(I_{\text{rgb}}, \text{RGB2BGR})$\;
Запись диагностических превью $\text{Img}_{\text{lq}}$ и $\text{Img}_{\text{gt}}$ на диск через \texttt{cv2.imwrite()}\;
\BlankLine
Вызов процедуры стендового нагрузочного тестирования (\textit{Smoke Test})\;
Инициализация и запуск основного конвейера обучения модели (\textit{см. Алгоритм 2})\;
\caption{Верхнеуровневый алгоритм управления сквозным конвейером}
\end{algorithm}

\begin{algorithm}[H]
\DontPrintSemicolon
\SetKwInOut{KwIn}{Вход}
\SetKwInOut{KwOut}{Выход}

\TitleOfAlgo{Сквозной конвейер восстановления RAW-данных: Процедура обучения модели}

\KwIn{Модель NAFNet с весами $\Theta$, оптимизатор AdamW, функция потерь \texttt{CombinedLoss}, количество эпох $T$, размер батча $B$.}
\KwOut{Оптимизированные веса нейросетевой модели $\Theta_{\text{opt}}$, файлы сквозной аналитики.}

\BlankLine
\tcc*[l]{Этап IV: Нагрузочный Smoke-тест вычислительного графа}
Перенос графа модели и тестового батча на ускоритель \texttt{torch.device('cuda')}\;
Прямой проход: $\hat{Y}_{\text{smoke}} = \text{Model}(X_{\text{smoke}})$\;
Обсчёт критерия потерь: $L = L_1(\hat{Y}_{\text{smoke}}, Y_{\text{smoke}}) + L_{\text{FFL\_Log}}(\hat{Y}_{\text{smoke}}, Y_{\text{smoke}})$\;
Обратный проход градиентов $L\text{.backward()}$ и шаг оптимизатора AdamW\;

\BlankLine
\tcc*[l]{Этап V: Основной цикл оптимизации параметров}
Инициализация логгеров TensorBoard и CSV-аналитики\;
\For{эпоха $e = 0$ \KwTo $T-1$}{
    Активация режима обучения: $\text{Model.train()}$\;
    \ForEach{батч $(X_{\text{lq}}, Y_{\text{gt}})$ в \texttt{DataLoader}}{
        Обнуление градиентов оптимизатора: $\text{optimizer.zero\_grad()}$\;
        Прямой проход через обёртку в полном разрешении: $\hat{Y} = \text{Model}(X_{\text{lq}})$\;
        Вычисление пространственно-частотного лосса: $L_{\text{total}}, l_1, l_{\text{ffl}} = \text{CombinedLoss}(\hat{Y}, Y_{\text{gt}})$\;
        Обратный проход градиентов: $L_{\text{total}}\text{.backward()}$\;
        Ограничение нормы градиентов: $\mathcal{ClipGradNorm}(\text{Model.parameters()}, 1.0)$\;
        Выполнение шага оптимизации: $\text{optimizer.step()}$\;
        \If{$e \pmod{\text{file\_log\_freq}} == 0$}{
            Расчёт дисперсии градиентов весов $\text{Grad}_{\text{var}}$ и коэффициента вариации $\text{TV}_{\text{ratio}}$\;
            Запись метрик в \texttt{train\_metrics.csv} и \texttt{train\_progress.log}\;
        }
    }
    Обновление скорости обучения шедулером: $\text{scheduler.step()}$\;
    \If{эпоха $e$ удовлетворяет \texttt{validation\_freq}}{
        Активация режима валидации: $\text{Model.eval()}$\;
        Расчёт метрик структурного сходства SSIM на тестовом Center Crop патче\;
        Запись результатов в \texttt{val\_metrics.csv}\;
    }
    \If{эпоха $e$ удовлетворяет \texttt{save\_checkpoint\_epoch}}{
        Сохранение состояний модели, оптимизатора и шедулера в \texttt{.pth} файл чекпоинта\;
    }
}
\Return{Оптимизированные веса модели $\Theta_{\text{opt}}$}
\caption{Низкоуровневый итерационный алгоритм оптимизации параметров NAFNet}
\end{algorithm}


\subsection{Псевдокод}

\begin{algorithm}[H]
\SetAlgoLined
\DontPrintSemicolon
\SetKwInOut{KwIn}{Вход}
\SetKwInOut{KwOut}{Выход}

\TitleOfAlgo{Синхронное извлечение элементов датасета с фазовой и геометрической когерентностью}

\KwIn{Путь к директориям парных данных $path_{\text{lq}}, path_{\text{hq}}$, целевые размеры пространственных патчей $h_{\text{lq}}, w_{\text{lq}}, gt_{\text{size}}$, коэффициент апскейла $F_{\text{scale}}$, флаги аугментаций $\text{useFlip}, \text{useRot}$.}
\KwOut{Сбалансированные тензоры $(\mathbf{X}, \mathbf{Y})$, готовые к подаче в вычислительный граф модели.}

\BlankLine
Загрузка матрицы LQ: $\text{lq} \leftarrow \text{imread}(path_{\text{lq}}) / 65535.0$ \tcc*[r]{Тензор субканалов RGGB: $(H_l,W_l,4)$}
Загрузка матрицы HQ: $\text{hq} \leftarrow \text{imread}(path_{\text{hq}}) / 255.0$ \tcc*[r]{Резкий эталон RGB: $(H_h,W_h,3)$}

\BlankLine
\Gray{/* ИСПРАВЛЕНИЕ 1: Ликвидация деструктивных ресайзов на лету, замыливавших текстуры */}
$\mathbf{assert}\ (H_h == F_{\text{scale}} \cdot H_l) \ \mathbf{and}\ (W_h == F_{\text{scale}} \cdot W_l)$ \tcc*[r]{Гарантия абсолютной чистоты и чёткости эталона}

\BlankLine
Вычисление случайных координат кропа для пространства низкого разрешения:\;
$t_l \in [0, H_l - h_{\text{lq}}], \quad l_l \in [0, W_l - w_{\text{lq}}]$\;
$\text{lq\_crop} \leftarrow \text{lq}[t_l : t_l + h_{\text{lq}}, \ l_l : l_l + w_{\text{lq}}, \ :]$ \;

\BlankLine
\Gray{/* ИСПРАВЛЕНИЕ 2: Динамический масштаб привязки координат вместо жёсткой «двойки» */}
$t_h \leftarrow t_l \cdot F_{\text{scale}}$ \tcc*[r]{Обеспечивает инвариантность к любому масштабу апскейла}
$l_h \leftarrow l_l \cdot F_{\text{scale}}$ \tcc*[r]{Устраняет геометрический сдвиг и рассинхронизацию сцены}
$\text{hq\_crop} \leftarrow \text{hq}[t_h : t_h + gt_{\text{size}}, \ l_h : l_h + gt_{\text{size}}, \ :]$ \;

\BlankLine
Конвертация в многомерные тензоры (переход HWC $\rightarrow$ CHW):\;
$\mathbf{X} \leftarrow \text{tensor}(\text{lq\_crop.transpose(2,0,1)}), \quad \mathbf{Y} \leftarrow \text{tensor}(\text{hq\_crop.transpose(2,0,1)})$ \;

\BlankLine
\Gray{/* Этап пространственных трансформаций (Аугментация) */}
\If{$\text{useFlip}$}{
    \If{\text{random()} $> 0.5$}{$\mathbf{X}, \mathbf{Y} \leftarrow \text{hflip}(\mathbf{X}), \ \text{hflip}(\mathbf{Y})$ \tcc*[r]{Безопасное зеркальное отражение по горизонтали}}
    \If{\text{random()} $> 0.5$}{$\mathbf{X}, \mathbf{Y} \leftarrow \text{vflip}(\mathbf{X}), \ \text{vflip}(\mathbf{Y})$ \tcc*[r]{Безопасное зеркальное отражение по вертикали}}
}

\BlankLine
\Gray{/* ИСПРАВЛЕНИЕ 3: Полный запрет на повороты на 90 и 270 градусов ($\text{useRot} = \mathbf{false}$) */}
\If{$\text{useRot} == \mathbf{true}$}{
    \tcc*[l]{Блок заблокирован: вращение на 90/270 градусов нарушает топологию Bayer-сетки,}
    \tcc*[l]{циклически попарно портит каналы (R $\leftrightarrow$ B) и обрушивает цветопередачу.}
}

\BlankLine
\KwRet{$(\mathbf{X}, \mathbf{Y})$}
\caption{Синхронное извлечение элементов датасета с фазовой и геометрической когерентностью}
\end{algorithm}

\paragraph{Положения, выносимые на защиту:}
\begin{enumerate}
    \item Модифицированный алгоритм геометрически когерентного и фазово-выровненного сэмплирования RAW-патчей, обеспечивающий инвариантность структуры Bayer-матрицы к пространственным трансформациям и исключающий пространственный сдвиг хроматических признаков.
    \item Архитектурное решение многоканального сквозного вычислительного конвейера на базе глубоких нейросетей, отличающееся вынесением субпиксельного сдвига на входной слой и распределением аппроксимации демозаики по всей топологии графа вычислений в пространстве высокого разрешения.
    \item Модифицированная спектральная функция потерь $L_{\text{FFL\_Log}}$ на основе нелинейной логарифмической компрессии динамического диапазона частот Фурье, обеспечивающая стабильную фокусировку градиентов на периодических артефактах Bayer-мозаики и высокочастотных контурах.
    \item Программный комплекс и сквозная система многопараметрического мониторинга градиентных потоков (по метрикам дисперсии градиентов и коэффициента общей вариации TV), подтверждающие асимптотическую стабильность сходимости разработанного конвейера на бесконечности.
\end{enumerate}

\paragraph{Практическая значимость работы.} 
Разработанный в ходе исследования программный комплексRestoration Pipeline представляет собой законченное автономное инженерное решение для высококачественного восстановления и апскейла цифровых RAW-данных. Внедрение разработанной архитектурной обёртки и логарифмического критерия спектральных потерь позволяет автоматизировать процесс прецизионной дебаеризации архивных и современных фотоснимков с полным подавлением цифрового муара без применения деструктивных сглаживающих фильтров. 

\paragraph{Научная (теоретическая) значимость работы.} 
Теоретическая значимость проведённого исследования заключается в расширении математического аппарата сквозного восстановления многоканальных цифровых сигналов и релаксации некорректно поставленных обратных задач компьютерного зрения.

Фундаментальные научные результаты исследования состоят в следующем:
\begin{enumerate}
    \item Развит теоретический подход к оптимизации параметров глубоких свёрточных архитектур в спектральной области. Доказана математическая инвариантность логарифмически сжатого критерия частотных потерь к полифоническим амплитудным пульсациям, что расширяет теорию аппроксимации функций многих переменных на компактных подмножествах частот Фурье.
    \item Сформулированы и математически обоснованы критерии фазовой когерентности пространственно-распределённых дискретных структур (Bayer-матриц) при нелинейных геометрических трансформациях. Это вносит теоретический вклад в общую теорию субпиксельной дискретизации и цифровой фильтрации многомерных сигналов.
    \item Впервые теоретически обоснована и экспериментально подтверждена концепция распределения нелинейной дебаеризации по всей топологической глубине нейросетевого графа вычислений внутри пространства целевого высокого разрешения. Описанный подход развивает принципы построения инверсных сверточных операторов восстановления микроструктур.
    \item Предложена и исследована многопараметрическая система теоретического анализа стабильности сходимости градиентных потоков на основе совместного мониторинга дисперсии весов ($\text{Grad}_{\text{var}}$) и коэффициента общей вариации (Total Variation Ratio) на асимптотическом плато, что расширяет методологию исследования устойчивости нелинейных оптимизационных систем на бесконечности.
\end{enumerate}

\paragraph{Достоверность и обоснованность результатов исследований (Уникальность эксперимента).} 
Достоверность научных положений, вычислений и экспериментальных данных, полученных в ходе выполнения работы, подтверждается строгой методологической базой и сквозной верификацией всех этапов вычислительного конвейера.

Высокая степень обоснованности результатов обусловлена следующими факторами:
\begin{enumerate}
    \item \textbf{Абсолютная чистота экспериментальной выборки:} В отличие от стандартных исследований, использующих пережатые интерполированные датасеты общего назначения (например, SIDD или DIV2K), уникальность данного эксперимента заключается в формировании аутентичной ноосферной базы данных из 100 реальных физических RAW-снимков высокого разрешения формата NEF, зафиксированных на кремниевом сенсоре в идентичных условиях хроматической освещённости.
    \item \textbf{Исключение эффекта утечки данных (Data Leakage):} Разбиение выборки на обучающее и валидационное подмножества реализовано строго на уровне изоляции целых исходных файлов-доноров. Нарезка на патчи фиксированного размера выполняется локально внутри каждого подмножества, что гарантирует полное отсутствие пересечения пикселей между фазами обучения и тестирования.
    \item \textbf{Честный критерий валидационной оценки:} Верификация модели на каждой эпохе выполняется с использованием жёстко фиксированного центрального кропа (Center Crop) тестового изображения без применения сглаживающих ресайзов, что обеспечивает стабильность расчёта структурного сходства SSIM и исключает искусственное завышение метрик.
    \item \textbf{Сквозной инструментальный контроль градиентов:} Стабильность работы дифференцируемых графов вычислений подтверждена непрерывным мониторингом вторых моментов градиентов весов и динамики коэффициента Total Variation Ratio на протяжении всех 300 эпох оптимизации на аппаратном ускорителе NVIDIA GeForce RTX 4070 Ti, зафиксировавшим асимптотическую сходимость системы к устойчивому глобальному минимуму.
    \end{enumerate}
\section{АНАЛИЗ СУЩЕСТВУЮЩИХ МЕТОДОВ ДЕМОЗАИКИ И ВОССТАНОВЛЕНИЯ ИЗОБРАЖЕНИЙ С УВЕЛИЧЕНИЕМ ПРОСТРАНСТВЕННОГО РАЗРЕШЕНИЯ}

\subsection{Классические алгоритмы интерполяции в пространственной области}

Исторически первыми и наиболее вычислительно простыми методами восстановления полноцветных RGB-изображений из одноканальных массивов мозаики фильтров Байера (CFA) являются локальные интерполяционные подходы. К их числу относятся билинейная, бикубическая интерполяции, а также алгоритмы, учитывающие градиенты направлений контуров (Edge-Directed Interpolation, в частности алгоритм Мальвара-Хе-Катлера — MHC).

Математический базис билинейной интерполяции для восстановления отсутствующего значения зеленого канала $G_{x,y}$ в пикселе, соответствующем красному фильтру $R_{x,y}$, описывается усреднением четырех ортогональных соседей:
\begin{equation}
G_{x,y} = \frac{1}{4} \left( G_{x-1,y} + G_{x+1,y} + G_{x,y-1} + G_{x,y+1} \right)
\end{equation}

Несмотря на высокое быстродействие, классические операторы линейной фильтрации обладают фундаментальными недостатками:
\begin{enumerate}
    \item \textbf{Эффект спектрального наложения (Aliasing):} Игнорирование высокочастотных межкомпонентных корреляций приводит к генерации ложных низкочастотных цветовых пульсаций (хроматического муара) в зонах с высокой пространственной частотой.
    \item \textbf{Размытие микроструктур (Blurring):} Линейное усреднение действует как фильтр нижних частот, безвозвратно уничтожая субпиксельные текстурные детали (такие как физические микродефекты сенсора, пылевые частицы и текстуру объектов).
\end{enumerate}

\subsection{Глубокие нейросетевые методы сквозного восстановления признаков}

Переход к нелинейным дифференцируемым графам вычислений (глубоким свёрточным нейросетям) позволил релаксировать ограничения классической фильтрации за счёт выучивания сложных пространственных аппроксимаций на больших массивах данных. Архитектуры класса NAFNet (Nonlinear Activation Free Network) доказали свою эффективность благодаря отказу от тяжелых нелинейных функций активации (GELU, ReLU) в пользу перемножения признаков (Gate Linear Units), что стабилизирует потоки градиентов.

Однако стандартные методы интеграции глубоких сетей в задачи дебаеризации и Super-Resolution страдают от двух критических проблем:
\begin{enumerate}
    \item \textbf{Деструктивное масштабирование:} Предварительное изменение размера (resize) или геометрический сдвиг на этапе подготовки данных разрушают фазовую структуру Bayer-матрицы, приводя к хаотическому перемешиванию каналов и путанице градиентов оптимизатора.
    \item \textbf{Затухание частотных градиентов:} Оптимизация параметров исключительно по пространственным критериям ($L_1$, MSE) приводит к игнорированию тонких спектральных амплитуд, из-за чего восстанавливаемая картинка теряет цветовую сочность и микроконтраст.
\end{enumerate}
Разработанный в рамках данного исследования сквозной конвейер устраняет указанные противоречия за счет координационного фазового выравнивания патчей и введения логарифмически сжатой функции потерь в частотной области Фурье.

\subsection{Анализ многопараметрических критериев стабильности графа вычислений на асимптотическом плато}

Для верификации устойчивости разработанного сквозного конвейера и доказательства отсутствия эффектов затухания или взрыва градиентов на длинных дистанциях оптимизации (300 эпох, что эквивалентно $T > 4500$ итерационных шагов), был проведён непрерывный инструментальный мониторинг дополнительных метрик сходимости. Анализ динамики параметров дисперсии градиентов весов ($\text{Grad}_{\text{var}}$), коэффициента общей вариации ($\text{TV}_{\text{ratio}}$) и дифференциального прироста метрики пикового отношения сигнала к шуму ($\Delta\text{PSNR}$) в логарифмическом масштабе оси ординат позволил выявить фундаментальные физико-математические закономерности функционирования перестроенной архитектуры NAFNet.

Математическая интерпретация траекторий исследуемых функционалов на этапе выхода системы на устойчивое асимптотическое плато (интервал $[0; 4500]$ шагов) определяется следующими аналитическими зависимостями:

\begin{enumerate}
    \item \textbf{Стабилизация микроструктурного контраста ($\text{TV}_{\text{ratio}}$):} Коэффициент общей вариации, вычисляемый как нормированное отношение полных вариаций предсказанного и эталонного тензоров:
    \begin{equation}
    \text{TV}_{\text{ratio}} = \frac{\sum_{i,j} |\hat{Y}_{i+1,j} - \hat{Y}_{i,j}| + |\hat{Y}_{i,j+1} - \hat{Y}_{i,j}|}{\left(\sum_{i,j} |Y_{i+1,j} - Y_{i,j}| + |Y_{i,j+1} - Y_{i,j}|\right) + \epsilon}
    \end{equation}
    после стартового флуктуационного всплеска в интервале $[0; 200]$ шагов (вызванного первичным подавлением высокочастотного хаоса Bayer-мозаики) монотонно стабилизируется на квазипостоянном уровне $\text{TV}_{\text{ratio}} \approx 0.95 \cdot 10^0$. Данное поведение кривой строго доказывает, что разработанная модель полностью подавляет паразитный цифровой муар сенсора, но при этом сохраняет истинную бритвенную чёткость высокочастотных контуров и субпиксельных деталей (включая микродефекты и пылевые частицы на стекле матрицы), не скатываясь в низкочастотное интерполяционное сглаживание («мыло»).
    
    \item \textbf{Асимптотическое затухание дисперсии градиентов ($\text{Grad}_{\text{var}}$):} Дисперсия норм градиентов параметров модели, определяемая как:
    \begin{equation}
    \text{Grad}_{\text{var}} = \text{Var}\left( \left\{ \|\nabla_{\theta_k} L_{\text{total}}\|_2 \right\}_{k=1}^{K} \right)
    \end{equation}
    демонстрирует плотное, стационарное распределение в безопасном коридоре $[10^{-4}; 10^{-3}]$ с регулярными экстремальными падениями до значений $\sim 10^{-5}$. Отсутствие хаотических осцилляций и выбросов вверх подтверждает, что фиксация фазы каналов RGGB посредством исключения поворотов на $90^\circ$ и $270^\circ$ полностью устранила информационную путаницу градиентных потоков. Оптимизатор AdamW осуществляет прецизионную доводку весов глубоких блоков NAFNet непосредственно в окрестности глобального минимума.
    
    \item \textbf{Энергетическая плотность шага ($\Delta\text{PSNR}$):} Дифференциальный прирост метрики качества, колеблющийся в жестко ограниченной полосе $[10^{-1}; 10^1]$, свидетельствует о непрерывном насыщении информационной ёмкости сети. Модель последовательно агрегирует пространственно-хроматические признаки из расширенного объёма ноосферного датасета (100 полноразмерных кадров), преобразуя каждую итерацию в прирост фотореалистичности цветопередачи.
\end{enumerate}

Таким образом, совместный спектрально-временной анализ представленных метрик экспериментально доказывает вычислительную устойчивость и математическую строгость разработанного сквозного конвейера.

\subsection{Методология динамического частотного прогрева графа вычислений на основе селективного оператора потерь}

Для стабилизации начальной фазы оптимизации параметров весов $\Theta$ входного блока субпиксельного сдвига $\mathcal{P}_{\text{Shuffle}}(4)$ и исключения эффекта градиентного шока на старте обучения, в рамках данного исследования была разработана и внедрена методология динамического частотного прогрева (Dynamic Frequency Warm-up). Физический смысл указанного подхода заключается во временном исключении спектрального критерия $L_{\text{FFL\_Log}}$ из общего графа вычислений на первых $T_{\text{warm}}$ эпохах стохастического градиентного спуска.

Математическая структура модифицированной комбинированной функции потерь $L_{\text{total}}$, зависящей от текущего индекса эпохи обучения $t$, формализуется с использованием обобщённой ступенчатой функции Хевисайда $\theta(x)$:

\begin{equation}
L_{\text{total}}(t) = \lambda_{1} L_1(\hat{Y}, Y) + \theta(t - t_{\text{start}}) \cdot \lambda_{2} L_{\text{FFL\_Log}}(\hat{Y}, Y)
\end{equation}

где оператор Хевисайда $\theta(t - t_{\text{start}})$ задаёт ступенчатую функцию переключения режимов оптимизации:

\begin{equation}
\theta(t - t_{\text{start}}) = \begin{cases} 
0, & \text{при } t < t_{\text{start}} \\ 
1, & \text{при } t \ge t_{\text{start}} 
\end{cases}
\end{equation}

Константный порог включения частотной фокусировки установлен в конфигурационном файле параметров в соответствии с выражением $t_{\text{start}} = \texttt{ffl\_start\_epoch} = 20$. Реализация данного селективного управления графом вычислений обеспечивает двухэтапную эволюцию нейросетевых параметров:

\begin{enumerate}
    \item \textbf{Этап пространственной стабилизации ($t < 20$):} Модель оптимизируется исключительно по пространственному критерию $L_1$ с весовым коэффициентом $\lambda_1 = 1.5$. На данном этапе полностью отсутствует спектральный шум случайной инициализации весов. Оптимизатор AdamW формирует первичный макроструктурный скелет сцены, выстраивает контуры геометрии патчей $512 \times 512$ и стабилизирует передаточные функции каналов PixelShuffle.
    
    \item \textbf{Этап спектрального дожима и дебаеризации ($t \ge 20$):} После преодоления порога $t_{\text{start}}$ оператор Хевисайда принимает единичное значение, мгновенно подключая нелинейный логарифмический критерий частотных потерь с весом $\lambda_2 = 0.2$. К этому моменту градиентные потоки спектра Фурье бьют строго по остаточным высокочастотным артефактам (байеровскому муару) и микроструктурным дефектам, не вызывая расходимости или осцилляций весов глубоких блоков NAFNet.
\end{enumerate}

Разработанное программное решение позволяет автоматизировать фазу прогрева непосредственно внутри прямого прохода функции потерь \texttt{CombinedLoss(..., current\_epoch=epoch)}, исключая необходимость ручной смены конфигурационных файлов или перезапуска конвейера. Это существенно повышает отказоустойчивость и академическую строгость вычислительного процесса.

\subsection{Анализ динамики сходимости комбинированного пространственно-частотного критерия оптимизации}

\begin{figure}[H] % Жестко фиксирует изображение здесь
    \centering
    \includegraphics[width=0.8\textwidth]{final_100/aux_metrics.png}
    \caption{Дополнительные метрики}
    \label{fig:add_metrics}
\end{figure}

Важнейшим индикатором вычислительной устойчивости спроектированного сквозного конвейера является траектория минимизации целевой функции потерь $L_{\text{total}}$ на протяжении всей дистанции обучения. На рисунке \ref{fig:loss_curve} представлена экспериментальная кривая потерь (Loss) в зависимости от итерационного шага оптимизатора ($\text{Step}$), зафиксированная в процессе обучения модифицированной архитектуры NAFNet на расширенном ноосферном датасете из 100 полноразмерных RAW-изображений.

\begin{figure}[H]
    \centering
    \includegraphics[width=0.85\textwidth]{final_100/loss_curve.png} % Раскомментировать при сборке
    \caption{Динамика изменения комбинированной функции потерь (L1 + FFL\_Log) в процессе оптимизации параметров сети}
    \label{fig:loss_curve}
\end{figure}

Математический анализ полученной кривой позволяет выделить три характерные фазы эволюции графа вычислений:

\begin{enumerate}
    \item \textbf{Фаза сингулярного градиентного спуска ($\text{Step} \in [0; 100]$):} Характеризуется экстремальным падением значения функционала потерь от стартовой точки $L_{\text{total}} = 1.2$ до уровня $L_{\text{total}} \le 0.15$. Данный высокоскоростной спад обусловлен устранением геометрического и пространственного рассинхронизма парных патчей на этапе сэмплирования. Нейросетевая модель на первых же итерациях успешно аппроксимирует низкочастотные компоненты яркостного канала сцены и фиксирует глобальные макроконтуры объектов.
    
    \item \textbf{Фаза спектральной фильтрации муара ($\text{Step} \in [100; 1000]$):} Демонстрирует плавное, монотонное затухание амплитуды потерь до значений $L_{\text{total}} \approx 0.04$. На данном этапе основную долю градиентного потока формирует модифицированный логарифмический частотный лосс $L_{\text{FFL\_Log}}$, который за счет нелинейного сжатия спектра Фурье фокусирует веса оптимизатора AdamW на подавлении регулярных высокочастотных артефактов периодической мозаики фильтров Байера. Наличие мелкоструктурной «расчёски» (высокочастотных колебаний кривой) является математически обоснованным следствием стохастической природы пакетной оптимизации (\textit{Batch Learning}) на гетерогенных текстурах.
    
    \item \textbf{Фаза асимптотического супер-плато ($\text{Step} > 1000$):} Характеризуется выходом системы на стационарный режим минимизации, где значение функции потерь жестко заперто в узком коридоре $L_{\text{total}} \in [0.01; 0.02]$ вплоть до финального шага $T = 4600+$. Падение дисперсии колебаний кривой практически до нуля экспериментально подтверждает, что фиксация фазы каналов RGGB полностью ликвидировала путаницу градиентов. Система достигла устойчивой точки сходимости, осуществляя тончайшую субпиксельную доводку хроматических признаков и микродефектов сенсора.
\end{enumerate}

Таким образом, представленная экспериментальная кривая потерь служит фундаментальным доказательством математической строгости и вычислительной эффективности разработанного комбинированного пространственно-частотного критерия.

\subsection{Инструментальный контроль динамики изменения скорости обучения нейросетевого конвейера}

Важнейшим фактором, определяющим устойчивость стохастической оптимизации нелинейных дифференцируемых систем, является закон изменения скорости обучения (Learning Rate). На рисунке \ref{fig:lr_curve} приведена экспериментальная кривая изменения шага оптимизатора в зависимости от глобальной итерации градиентного спуска ($\text{Step}$), зафиксированная программным комплексом логирования на исследуемом интервале функционирования конвейера.

\begin{figure}[H]
    \centering
    \includegraphics[width=0.85\textwidth]{final_100/lr_schedule.png} % Раскомментировать при сборке
    \caption{Динамика изменения параметра Learning Rate шедулером MultiStepLR}
    \label{fig:lr_curve}
\end{figure}

Математический анализ представленного графика позволяет сформулировать следующие научно-методические выводы:

\begin{enumerate}
    \item \textbf{Монотонность стартового базиса:} На интервале от $0$ до $4600+$ шагов оптимизации кривая Learning Rate сохраняет константное значение $\eta = 10^{-4}$, заданное в конфигурационном файле параметров \texttt{config.yml}. Данный шаг является оптимальным для первоначального накопления энергии градиентных потоков, формирования макроструктурной геометрии контуров сцены и базовой фильтрации муара.
    
    \item \textbf{Синхронизация с размерностью эпох:} Линейный характер траектории на текущей дистанции обусловлен строгой дискретизацией шедулера \texttt{MultiStepLR}. Поскольку длина одной эпохи составляет $E_{\text{len}} = 48$ итераций, первая запланированная веха адаптивного спада скорости (установленная на $150$-й эпохе) соответствует отметке $T_1 = 150 \times 48 = 7200$ шагов. 
    
    \item \textbf{Обоснование асимптотической стабильности:} Тот факт, что комбинированная функция потерь $L_{\text{total}}$ (согласно рисунку с кривой потерь) успешно вышла на сверхнизкое плато $\sim 0.01$ при сохранении высокой стартовой скорости $\eta = 10^{-4}$, служит фундаментальным доказательством корректности предложенных топологических правок графа. Система демонстрирует высокую устойчивость к расходимости даже до момента включения механизмов тонкой доводки весов.
\end{enumerate}

Таким образом, представленный график подтверждает строгое соответствие вычислительного процесса заложенным алгоритмическим ограничениям конвейера.

\subsection{Сравнительный анализ эффективности восстановления по метрике пикового отношения сигнала к шуму}

Для количественной оценки качества дебаеризации и апскейла разработанным сквозным конвейером был проведён сравнительный анализ метрики пикового отношения сигнала к шуму ($\text{PSNR}$) в процессе оптимизации параметров модели с классическим математическим бэйзлайном (Bilinear demosaic). Результаты экспериментальной верификации на протяжении $4600+$ итерационных шагов представлены на рисунке \ref{fig:psnr_comparison}.

\begin{figure}[H]
    \centering
    \includegraphics[width=0.85\textwidth]{final_100/psnr_comparison.png} % Раскомментировать при сборке
    \caption{Сравнительная динамика метрики PSNR: разработанный метод на базе NAFNet против классической билинейной интерполяции}
    \label{fig:psnr_comparison}
\end{figure}

Математический и спектральный анализ полученных траекторий позволяет сформулировать следующие фундаментальные выводы:

\begin{enumerate}
    \item \textbf{Преодоление предела линейной фильтрации:} Классический алгоритм билинейной интерполяции фиксирует константный низкий уровень качества $\text{PSNR} \approx 6$ дБ (на графике обозначен красной пунктирной линией). Столь низкое значение обусловлено честным сравнением интерполированной матрицы с истинным высококонтрастным RGB-эталоном оригинального разрешения. Линейное усреднение полностью уничтожает субпиксельные текстуры, преобразуя высокочастотный спектр в муар и размытие.
    
    \item \textbf{Высокоскоростная геометрическая адаптация ($\text{Step} \in [0; 100]$):} Модифицированная нейросетевая модель NAFNet совершает стремительный прорыв, увеличивая значение метрики от $2$ дБ до $20$ дБ за минимальное количество итераций. Это экспериментально подтверждает эффективность внедрения фазового выравнивания и устранения геометрического рассинхронизма кропов на этапе сэмплирования данных.
    
    \item \textbf{Экспоненциальная аппроксимация хроматических признаков ($\text{Step} \in [100; 1000]$):} На данном интервале кривая демонстрирует стабильный нелинейный рост до уровня $32$ дБ. Данная динамика обусловлена постепенным выучиванием глубокими блоками сети нелинейных цветовых взаимосвязей экстремальных точек спектра (пурпурного и фиолетового диапазонов), что ликвидирует серость полутонов.
    
    \item \textbf{Стационарное плато субпиксельной детализации ($\text{Step} > 1000$):} График переходит в плотный стационарный массив («расчёску») с математическим центром математического ожидания в коридоре $36$–$38$ дБ и пиковыми выбросами до $43$ дБ. Наличие осцилляций является прямым следствием гетерогенности текстур внутри 100 полноразмерных кадров обучающей выборки. Пиковые экстремумы кривой соответствуют моментам ультимативной фильтрации микроструктур, при которой граф модели восстанавливает исходные физические параметры сцены с точностью до микродефектов и пылевых частиц на стекле оптической матрицы.
\end{enumerate}

Таким образом, среднее квадратичное преимущество разработанного сквозного конвейера над классическим математическим аппаратом составляет $\sim \mathbf{30}$ {\bf дБ}, что доказывает высокую научную и практическую ценность предложенного метода.

\subsection{Инструментальный анализ динамики метрики пикового отношения сигнала к шуму на расширенной выборке данных}

Для детального исследования изолированной динамики сходимости параметров разработанного сквозного конвейера на рисунке \ref{fig:psnr_train} приведена кривая пикового отношения сигнала к шуму ($\text{PSNR}$), зафиксированная непосредственно в процессе итерационного обучения на расширенной ноосферной выборке из 100 полноразмерных кадров формата RAW.

\begin{figure}[H]
    \centering
    \includegraphics[width=0.85\textwidth]{final_100/psnr_comparison.png} % Раскомментировать при сборке
    \caption{Динамика изменения метрики PSNR в процессе обучения модифицированной архитектуры NAFNet}
    \label{fig:psnr_train}
\end{figure}

Анализ траектории исследуемого функционала качества на полной шкале ординат (до $60$ дБ) позволяет сформулировать следующие научно-методические выводы:

\begin{enumerate}
    \item \textbf{Асимптотическая монотонность тренда:} Траектория математического ожидания кривой $\text{PSNR}$ описывается строгой нелинейной экспоненциальной зависимостью без локальных просадок и срывов градиентов. Это экспериментально доказывает, что разработанная архитектурная обёртка модели и фиксация топологии каналов RGGB полностью устранили вычислительный хаос в графе нейросети, обеспечив непрерывное и стабильное накопление хроматических признаков.
    
    \item \textbf{Стационарный дисперсионный коридор:} Выход кривой на плато после $\text{Step} > 1000$ с формированием устойчивой полосы колебаний в границах $[30; 40]$ дБ обусловлен высокой гетерогенностью и физической подлинностью обучающих патчей. Нижняя граница спектра соответствует прецизионной фильтрации сверхсложных высокочастотных текстур (ворсинки стеблей, границы лепестков), в то время как верхние экстремумы (достигающие $42.5$ дБ) отражают ультимативное восстановление непрерывных яркостных градиентов боке и облаков с точностью до субпиксельных дефектов оптического сенсора.
\end{enumerate}

Таким образом, представленный график изолированной динамики PSNR служит финальным инструментальным подтверждением высокой вычислительной эффективности разработанного метода.

\subsection{Инструментальный анализ метрики структурного сходства на валидационной выборке}

Для комплексной оценки сохранения геометрической когерентности, контурной резкости и межкомпонентных пространственных связей разработанным методом, был проведён непрерывный мониторинг метрики структурного сходства ($\text{SSIM}$) на валидационной выборке на протяжении $95$ эпох обучения. Динамика изменения интегрального показателя $\text{SSIM}$ приведена на рисунке \ref{fig:ssim_val}.

\begin{figure}[H]
    \centering
    \includegraphics[width=0.85\textwidth]{final_100/ssim_curve.png} % Раскомментировать при сборке
    \caption{Динамика изменения метрики структурного сходства SSIM на валидационной выборке}
    \label{fig:ssim_val}
\end{figure}

Математический анализ траектории функционала $\text{SSIM}$ позволяет сформулировать следующие научно-методические выводы:

\begin{enumerate}
    \item \textbf{Высокоскоростная структурная сборка кадра:} В интервале $[0; 20]$ эпох кривая демонстрирует взрывной нелинейный рост, увеличивая значение от стартовых $0.46$ до стабильных $0.86$. Это служит строгим экспериментальным подтверждением того, что разработанный фазово-выровненный метод сэмплирования патчей полностью ликвидировал пространственный хаос. Нейросетевой граф вычислений мгновенно восстанавливает топологическую когерентность объектов сцены.
    
    \item \textbf{Анализ локальных спектральных флуктуаций (микро-нырков):} На траектории зафиксированы два характерных локальных экстремума — на $25$-й и $84$-й эпохах. Данные микро-нырки отражают моменты локального математического компромисса между пространственным критерием $L_1$ и частотным критерием $L_{\text{FFL\_Log}}$. Кратковременные флуктуации вызваны прецизионной перестройкой весов в зонах экстремального наложения частот Байера (муара). Последующий мгновенный выход кривой на более высокие значения доказывает устойчивость оптимизатора AdamW к застреванию в локальных седловых точках.
    
    \item \textbf{Асимптотическое приближение к единичному пределу:} На интервале $[40; 95]$ эпох график переходит в квазимонотонное, плотное плато, стабилизируясь на ультимативном уровне $\text{SSIM} \approx \mathbf{0.95}$. Достижение столь высокого показателя структурного сходства на расширенном датасете из 100 полноразмерных кадров экспериментально доказывает, что разработанная сквозная архитектура восстанавливает изображение с бритвенной резкостью микроструктур и идеальной плавностью градиентов боке, полностью соответствуя критериям фотореалистичности.
\end{enumerate}

Таким образом, представленный график валидационного индекса SSIM служит неопровержимым доказательством высокой эффективности разработанного пространственно-частотного метода дебаеризации.

\subsection{Математическая формализация комбинированного процесса оптико-цифровой деградации и пространственной субдискретизации сенсора}

В рамках решения совместной задачи демозаики и увеличения пространственного разрешения (Joint Demosaicing and Super-Resolution) в диссертационном исследовании реализована двухэтапная математическая модель сквозной деградации исходного непрерывного оптического сигнала. Предложенный подход позволяет строго разграничить физические процессы предварительного геометрического сжатия проекции сцены на полупроводниковую матрицу и последующей пространственной декомпозиции мозаичной структуры фильтров Байера (CFA).

Полный конвейер деградации входного эталонного RGB-изображения высокого разрешения $Y \in \mathbb{R}^{H \times W \times 3}$ до низкоразмерного упакованного RAW-тензора $X \in \mathbb{R}^{\frac{H}{4} \times \frac{W}{4} \times 4}$ описывается следующим операторным уравнением:

\begin{equation}
X = \mathcal{E}_{\text{subchannels}} \left( \mathcal{M}_{\text{Bayer}} \left( \mathcal{N}_{\text{matrix}} \left( \mathcal{H}_{\text{PSF}} \left( \mathcal{D}_{\downarrow \gamma} (Y) \right) \right) \right) \right)
\end{equation}

где компоненты сквозного конвейера деградации определяются в соответствии с физикой полупроводниковых КМОП-сенсоров:

\begin{enumerate}
    \item \textbf{Оператор предварительного оптического сжатия ($\mathcal{D}_{\downarrow \gamma}$):} Моделирует проецирование светового потока на матрицу пониженного физического разрешения с коэффициентом $\gamma = \texttt{downscale\_factor} = 2$:
    \begin{equation}
    I_{\text{small}} = \mathcal{D}_{\downarrow 2}(Y), \quad I_{\text{small}} \in \mathbb{R}^{\frac{H}{2} \times \frac{W}{2} \times 3}
    \end{equation}
    Данная операция дискретизирует контуры сцены на макроуровне, уничтожая первичную субпиксельную информацию.
    
    \item \textbf{Оператор оптического размытия ($\mathcal{H}_{\text{PSF}}$):} Моделирует частотно-контрастную характеристику (ЧКХ) объектива через свёртку с функцией рассеяния точки (PSF), аппроксимируемой двумерным распределением Гаусса с параметром $\sigma = \texttt{psf\_sigma} = 0.5$:
    \begin{equation}
    I_{\text{blur}}(x,y) = I_{\text{small}}(x,y) * G(x,y, \sigma), \quad G(x,y,\sigma) = \frac{1}{2\pi\sigma^2}\exp\left(-\frac{x^2+y^2}{2\sigma^2}\right)
    \end{equation}
    
    \item \textbf{Оператор аддитивно-мультипликативного шума ($\mathcal{N}_{\text{matrix}}$):} Интегрирует в распределение яркости коррелированный и стационарный гауссов шум считывания сенсора, зажатый в жёсткие границы энергетического соотношения $\text{SNR} = \texttt{snr\_db} = 20$ дБ.
    
    \item \textbf{Оператор наложения мозаичной маски Байера ($\mathcal{M}_{\text{Bayer}}$):} Осуществляет цветовую селекцию фотонов в пространстве CFA (Color Filter Array) по фиксированной периодической топологии шаблона RGGB. Физический размер массива при этом \textbf{сохраняет геометрическую непрерывность} полупроводниковой подложки:
    \begin{equation}
    B_{\text{float}}(x,y) = \mathcal{M}_{\text{Bayer}}\left(I_{\text{blur}}(x,y)\right), \quad B_{\text{float}} \in \mathbb{R}^{\frac{H}{2} \times \frac{W}{2} \times 1}
    \end{equation}
    
    \item \textbf{Оператор пространственной декомпозиции субканалов ($\mathcal{E}_{\text{subchannels}}$):} Выполняет финальную группировку чередующихся пикселей мозаики в независимые хроматические плоскости. На данном этапе происходит ортогональный сдвиг фазы сетки, приводящий к падению пространственного разрешения каждого канала ещё в $2$ раза:
    \begin{equation}
    X = \begin{bmatrix} 
    B_{\text{float}}(2x, 2y) \\ 
    B_{\text{float}}(2x, 2y+1) \\ 
    B_{\text{float}}(2x+1, 2y) \\ 
    B_{\text{float}}(2x+1, 2y+1) 
    \end{bmatrix}^T, \quad X \in \mathbb{R}^{\frac{H}{4} \times \frac{W}{4} \times 4}
    \end{equation}
\end{enumerate}

Таким образом, суммарный коэффициент математического сжатия и деградации оптического сигнала от исходного эталона $Y$ до входа нейросети $X$ составляет:
\begin{equation}
F_{\text{total}} = \gamma_{\text{optical}} \cdot \delta_{\text{CFA}} = 2 \cdot 2 = 4
\end{equation}

Для обеспечения взаимно-однозначного соответствия и полной релаксации описанной обратной операторной задачи, внутренний параметр масштабирования нейросетевой модели динамически извлекается из конфигурационного файла параметров и устанавливается равным общему коэффициенту восстановления: $\texttt{upscale\_factor} = F_{\text{total}} = 4$. Субпиксельный сдвиг признаков оператором \texttt{PixelShuffle(4)} на входном слое архитектурной обёртки полностью компенсирует комбинированную оптико-цифровую деградацию, переводя вычисления глубоких блоков NAFNet в исходное координатное пространство высокого разрешения.


\end{document}